\documentclass[letterpaper,spanish,12pt,openany]{book}
\usepackage[utf8]{inputenc}
\usepackage[spanish,es-tabla,mexico]{babel}
\usepackage[T1]{fontenc}
\usepackage{csquotes}
\usepackage{graphicx}
\usepackage{wrapfig}
\usepackage{verbatim}
\usepackage{amsmath, amssymb}
\usepackage{MnSymbol}
\usepackage{upgreek}
\usepackage{mathrsfs}
\usepackage[mathcal]{euscript}
%\usepackage{physics}
\usepackage{multicol}
\usepackage{multirow}
\usepackage[hidelinks]{hyperref}
\usepackage{endnotes}
\usepackage{array}
\usepackage{xcolor}
\usepackage{colortbl}
\usepackage{lscape}
\usepackage{enumitem}
\usepackage{units}
\usepackage{gensymb}
\usepackage{glossaries}
\usepackage{textgreek}
\usepackage{import}
\usepackage{caption}
\usepackage{lipsum}
\usepackage{subfig}
\usepackage{float}
\usepackage{commath}
\usepackage{amsbsy}
\usepackage[font=small]{caption}
\usepackage[left=2.5cm,right=2.5cm,top=2 cm,bottom=2 cm]{geometry}
\usepackage[backend=biber,sorting=none,style=numeric-comp,bibstyle=chem-acs,articletitle=true,biblabel=brackets,doi=false, maxcitenames=2, autocite = superscript]{biblatex} %Nota del 05 de julio de 2021: esta librería genera las referencias bibliográficas.
\bibliography{bib.bib}
\usepackage{verbatim} %Nota del 05 de julio de 2021: esta librería es para comentar varias lineas
\usepackage{float} %Nota del 06 de julio de 2021: esta librería sirve para las imágenes en el entorno multicols
\usepackage{ragged2e}
\usepackage[version=4]{mhchem}
\newcommand{\e}{\hat{\textbf{e}}} %para los vectores unitarios
\newcommand{\im}{\mathrm{i}} %para las partes imaginarias, es la i
\newcommand{\neod}{\ce{NdFeO3}}
\newcommand{\sama}{\ce{SmFeO3}}
\newcommand{\grado}{$^\circ$}
\newcommand{\gradoC}{$^\circ$C}
\usepackage{subfiles}
\usepackage{pdfpages}
\usepackage{titlesec}
\setcounter{secnumdepth}{4}
\titleformat{\paragraph}
{\normalfont\normalsize\bfseries}{\theparagraph}{1em}{}
\titlespacing*{\paragraph}
{0pt}{3.25ex plus 1ex minus .2ex}{1.5ex plus .2ex}
\usepackage{titlesec}
\titleformat{\chapter} % Config titlesec
	{\Large\bfseries}	% 
	{\huge \thechapter}	% 
	{20pt}				% 
	{\huge}				% 
\usepackage{mathtools}
	%\DeclarePairedDelimiter\bra{\langle}{\rvert} % \bra{a} -> <a|
	%\DeclarePairedDelimiter\ket{\lvert}{\rangle} % \ket{a} -> |a>
	%\DeclarePairedDelimiterX\braket[2]{\langle}{\rangle}{#1\,\delimsize\vert\,\mathopen{}#2} % \braket{a}{b} -> <a|b>
	%\DeclarePairedDelimiter\avrg{\langle}{\rangle} % \avrg{x} -> <x>
	\newcommand{\bra}[1]{\ensuremath{\langle #1 \rvert}}
	\newcommand{\ket}[1]{\ensuremath{\lvert #1 \rangle}}
	\newcommand{\braket}[2]{\ensuremath{\langle #1 \mid #2 \rangle}}
	\newcommand{\avrg}[1]{\ensuremath{\langle #1 \rangle}}
\begin{document}
%COMIENZA SECCIÓN COMENTADA
%-------------------------------------------
%\frontmatter
%\includepdf{0. Portada/0. Portada.pdf}
%\clearpage
%------------------------------------------
% DEDICATORIA
%------------------------------------------
%\chapter*{}
%\vspace{2cm}
%\begin{flushright}
%	{\fontfamily{ppl}%Palatino
%	\selectfont\emph{<<Cita. AAAAAAAAAAAAAAAAAAAAAAAAAAAAAAAAAAAAAAAAAAAAAAA AAAAAAAAAAAAAAAAAAAAAAAAAAAAAAAAAAAAAAAAAAAAAAA AAAAAAAAAAAAAAAAAAAAAAAAAAAAAAAAAAA.>>}\\
%\vspace{5mm}
%	Autor}

%\end{flushright}
%\vfill
%------------------------------------------
% AGRADECIMIENTOS
%------------------------------------------
%\chapter{Agradecimientos}
%\lipsum[6-9]
%-------------------------------------------------
%FIN DE SECCIÓN COMENTADA
\tableofcontents % Index 
\mainmatter 
%\subfile{1. Resumen/1. Resumen.tex}
%\subfile{2. Introduccion/2. Introduccion.tex}
\subfile{3. Estado del Arte/3. Estado del Arte.tex}
\subfile{4. Metodologia/4. Metodologia.tex}
\subfile{5. Resultados/5. Resultados.tex}
%\subfile{6. Conclusion/6. Conclusion.tex}
\printbibliography[heading=bibintoc, title={Bibliografía}]
%\subfile{8. Anexos/8. Anexos.tex}
\end{document}